\section{Research Questions and Falsifiable Outcomes}
\subsection{Primary question}
The publication program asks whether conventional indicators of successful forgetting
agree with behavioral consistency relative to models that never encountered the deleted
records.  Let $C_j(M_u)$ denote the pass/fail decision induced by conventional criterion
$j$, and let $E(M_u)$ denote consistency with the exact-retraining reference population.
The central estimand is not merely mean method performance, but the frequency and form of
the contradiction
\begin{equation}
  C_1(M_u)=\cdots=C_q(M_u)=\mathrm{pass},\qquad
  E(M_u)=\mathrm{fail}.
\end{equation}
This definition deliberately permits the criteria to disagree.  It does not presume that
counterfactual fidelity is one-dimensional or that every discrepancy is practically
important.

\subsection{Research questions}
The confirmatory program is organized around five questions.  \textbf{RQ1} asks whether a
deletion intervention is identifiable against exact-training randomness for the chosen
regime.  \textbf{RQ2} asks how often utility, forget-set, privacy, and calibration criteria
agree with exact-counterfactual consistency.  \textbf{RQ3} asks which fidelity dimensions
detect contradictions that aggregate accuracy misses.  \textbf{RQ4} asks whether repeated
deletion requests produce cumulative discrepancy.  \textbf{RQ5} asks how fidelity changes
as computational cost is reduced relative to retraining.  P001 addresses only the
measurement precondition in RQ1; the other questions remain prospective.

\subsection{Outcome taxonomy}
Each evaluated condition has four scientifically distinct outcomes.  A \emph{consistent}
outcome places the approximate model within the prespecified exact envelope on the primary
endpoint.  A \emph{contradictory} outcome passes the conventional criteria but fails that
envelope.  A \emph{concordant failure} fails both.  An \emph{uninformative} outcome occurs
when reference sensitivity, sample size, run integrity, or precision is inadequate for a
decision.  The last state is essential: absence of detectable discrepancy is not evidence
of equivalence when the assay cannot resolve the intervention.

\subsection{Falsification criteria}
The motivating measurement hypothesis would be weakened if conventional criteria and
counterfactual decisions agree across preregistered datasets, deletion regimes, methods,
and seeds; if apparent contradictions vanish under independently chosen equivalence
margins; or if counterfactual distances add no stable information beyond inexpensive
metrics.  Conversely, isolated exploratory contradictions are insufficient.  A defensible
claim requires repeated, integrity-valid observations under a frozen analysis plan and a
reference population demonstrably sensitive to the deletion regime.
